% Section 6 -- Results (merged Results + Analysis). Flowing moderate sentences; long sentences split.
% Numbers from results.json. Floats: T1, F3, F4, T4.
\section{Results}
\label{sec:results}

Reinforcement learning against the conditional-rate reward calibrates the forecaster to the level of
the betting market. On the held-out 2024 season the direct model reaches an ECE of
$0.029$, against the market's $0.027$ and the untrained model's $0.057$. Its maximum calibration error,
$0.060$, is below the market's $0.082$, and its reliability curve lies on the diagonal where the
untrained model's does not (\cref{fig:reliability}). Its Brier score of $0.144$ improves on the untrained
$0.206$ and trails the market's $0.136$ (\cref{tab:main}). The masked chain-of-thought model is as well
calibrated, at an ECE of $0.030$, with a slightly higher Brier score of $0.152$.

\begin{table*}[t]
\centering
% requires \booktabs. Brier shows 95% bootstrap CI (10k reps over plays).
\begin{tabular}{lccccc}
\toprule
System & Brier & ECE & MCE & Acc. & Resolution \\
\midrule
Qwen2.5-7B zero-shot (direct) & 0.2057\,{\scriptsize[0.2016,\,0.2098]} & 0.0569 & 0.2539 & 0.644 & 0.0388 \\
Qwen2.5-7B zero-shot (CoT) & 0.1681\,{\scriptsize[0.1626,\,0.1736]} & 0.0687 & 0.1493 & 0.756 & 0.0875 \\
Masked-CoT RLVR (ours) & 0.1522\,{\scriptsize[0.1466,\,0.1577]} & 0.0293 & 0.0684 & 0.777 & 0.0986 \\
Direct RLVR (ours) & 0.1443\,{\scriptsize[0.1394,\,0.1491]} & 0.0292 & 0.0596 & 0.784 & 0.1058 \\
DeepSeek-V4 zero-shot & 0.1438\,{\scriptsize[0.1392,\,0.1483]} & 0.0430 & 0.0716 & 0.790 & 0.1078 \\
\midrule
Empirical rate $\hat{p}$ & 0.1432\,{\scriptsize[0.1384,\,0.1480]} & 0.0437 & 0.0993 & 0.783 & 0.1083 \\
Betting market & 0.1355\,{\scriptsize[0.1307,\,0.1403]} & 0.0273 & 0.0824 & 0.799 & 0.1148 \\
\bottomrule
\end{tabular}

\caption{Held-out 2024 ($n=5185$). Brier with $95\%$ bootstrap interval, expected and maximum
calibration error, accuracy, and Murphy resolution. The trained models reach the teacher and
frontier-model tier, the direct model matches the market on calibration, and both trail the market on
Brier through resolution rather than reliability.}
\label{tab:main}
\end{table*}

\begin{figure*}[t]
\centering
\includegraphics[width=\linewidth]{F3_reliability.pdf}
\caption{Reliability on held-out 2024; points on the diagonal are perfectly calibrated and their area is
proportional to the plays in the bin. The trained forecasters follow the diagonal as closely as the
market, while the untrained chain-of-thought model is overconfident.}
\label{fig:reliability}
\end{figure*}

The direct model matches the market on calibration but trails it on Brier by $0.008$. By the Murphy
decomposition this gap is resolution, not reliability: the model is as well calibrated as the market but
less sharp, $0.106$ against $0.115$. It assigns the right probability to the states it can describe, but
separates fewer of them. The missing sharpness comes from information the prompt does not contain, and
no static forecaster does better.

The frontier model and the teacher reach the same Brier score as the direct model, $0.144$ and $0.143$
against $0.144$, and a paired bootstrap separates none of the three. All three trail the market by
$0.008$, significant in each case (\cref{fig:triangulation}), and the direct model is better calibrated
than the frontier model, $0.029$ against $0.043$. No richer static model improves on the coarse rate. The
tuned nflverse model and a gradient-boosted model on the full feature set score $0.156$ and $0.158$
against $0.143$ (\cref{tab:ceiling}), since the extra features add play-level noise without game-level
signal. \citet{brill2024} reach the same bound analytically, showing that the correlated structure of
play-by-play data limits the accuracy of any such model. The market's $0.008$ edge is therefore live
in-game information that no static input contains.

\begin{figure*}[t]
\centering
\includegraphics[width=\linewidth]{F4_triangulation.pdf}
\caption{Held-out 2024 Brier with $95\%$ bootstrap intervals. A label-free 7B model trained by direct
prediction, a frontier model prompted zero-shot, and the empirical-rate teacher converge near $0.143$
and trail the market by the same significant margin, placing the residual gap in the available
information rather than in model capacity.}
\label{fig:triangulation}
\end{figure*}

\begin{table}[t]
\centering
% static-feature Brier ceiling. Only the coarse buckets and the LLM reach the LLM tier;
% richer-feature static models are worse. Vegas's edge is live information.
\begin{tabular}{lcc}
\toprule
Static-feature forecaster & Brier & ECE \\
\midrule
Coarse $\hat{p}$ (our reward) & 0.1432 & 0.0437 \\
nflverse WP model & 0.1562 & 0.0188 \\
GBM, all features & 0.1584 & 0.0260 \\
\midrule
Direct RLVR (ours) & 0.1443 & 0.0292 \\
Betting market & 0.1355 & 0.0273 \\
\bottomrule
\end{tabular}

\caption{Held-out 2024 Brier and ECE for static-feature forecasters. Two models with more features than
the coarse buckets are worse on Brier; the coarse empirical rate is the best static estimator, and the
direct model reaches its level.}
\label{tab:ceiling}
\end{table}

On the test season the direct model is better calibrated than the teacher it was trained on. Its
ECE is $0.029$ against the teacher's $0.044$, and its maximum calibration error $0.060$
against $0.099$, at the same Brier score. The policy smooths the bucketed target into a calibrated
function of the state.

The decalibration of full-completion chain-of-thought training (\cref{sec:decoupling}) does not come
from misreading the game state. A blinded judge scored five state-reading errors (possession, score,
spread, clock, and pregame anchoring) and found each in one to ten percent of completions across the
base, full-completion, and masked models, so all of them read the game correctly. The full-completion
models fail instead at the step from those facts to a number. They write pseudo-quantitative arguments whose probabilities exceed one and
report a figure their own text does not support (\cref{fig:reasoning-examples}). The error is an
inconsistency between the reasoning and the answer (\cref{fig:faithfulness}), not a factual mistake.
Masking the gradient leaves the reading of the state untouched and keeps the number consistent with it,
which restores calibration.
